% Transcription — fonds Alexandre Grothendieck, shelfmark 2, batch 1 (pages 1-20).
% Established from the facsimile archives/batches/2/batch-01.pdf (a mirror of
% https://grothendieck.umontpellier.fr/2.pdf), per the skill
% transcribe-grothendieck. The batch file carries no cover sheet: PDF page N is
% archivists' page N. Checked against the pencilled numbers bottom-left ("1",
% "2", "4", "8", "13" read on the corresponding leaves).
%
% Pass: Opus 5 (claude-opus-5), 2026-09-19 — first pass, unchecked against the pages by a human.
%
% What the batch is. Three runs in his hand, each opened by a wrapper he
% inscribed, on loose leaves and on the versos of a typescript that is not his.
% The folder's inventory title is "Théorème de Riemann-Roch / RR Notes
% Antiques", and the first run is exactly that: on a non-singular curve, the
% ring E of classes of coherent sheaves modulo the short exact sequences, the
% subgroups E_0, N, D and the subalgebra Phi of vector-bundle classes, the
% identification E/N ≃ Z + P, and Riemann-Roch written as a single line,
% \int(F) = K_1( c(F)(1 - K/2) ). Pages 4, 6 and 7 restate the same thing for a
% noetherian preschema of dim ≤ 1, in the two-graded notation K^.(X) (bundles)
% and K_.(X) (coherent sheaves), with the gamma-filtration, Pic(X), pi_0(X) and
% the comparison map c = ch into B^.(X), B_.(X). Pages 9, 11 and 13 are the
% "Contre-exemples" run: the Chern class of an i_! along X -> X x P^r, computed
% through K(X x P^r) = K(X) ⊗ K(P^r)/(xi = 1 - L^{-1}), giving
% C~(i_!(x)) = C~(x) * (1 + (-1)^{r-1}(r-1)! eta^r), the comparison with R.R.,
% and on page 13 the elliptic curve for which the class comes out trivial.
% Pages 15 to 18 are a different subject altogether — the category A(X, G) of
% Modules on an annulated topos carrying an action of a group G, the
% restriction and extension functors, the exterior product ⊠, the product
% F * G over the symmetric groups, and the formulas for the T^n. Page 20 opens
% the lambda-ring formalism: a special lambda-ring with a family of generators,
% the universal ring B, the Chern ring CB, the gamma-filtration and the two
% maps psi_K and phi_K between the graded ring and CK.
%
% Skipped, without description, as the skill requires: page 5 (blank), pages
% 10, 12 and 14 (a typescript that is not his; page 14 carries one word of his
% correcting the typing, which is not mathematics in hand). Pages 1, 8 and 19
% are wrappers inscribed by him; their words are taken up as the section
% headings below and they get no \page, so that the gaps in the numbering
% record what was passed over.
%
% Notation fixed once here, and held for the batch. His two K's are
% K^{\bullet}(X) for vector bundles and K_{\bullet}(X) for coherent sheaves
% (he writes K with a raised and a lowered dot); likewise B^{\bullet},
% B_{\bullet}. The looped functional of pages 2 and 3, drawn as a long s, is
% transcribed \int. The looped and barred capital of pages 2 and 3, the group
% the divisors map onto, is transcribed \mathfrak{P}; the plain script L of
% page 2 is \mathcal{L} and the script D is \mathcal{D} — the three are
% distinct letters on the leaf. His capital E is drawn like a reversed 3 and is
% transcribed E. The tilde'd Chern class of pages 9-13 is \widetilde{C} and the
% ring it lands in \widetilde{A}; the hyperplane class of P^r in those pages is
% \eta. Underlined sheaves are \underline{\mathcal{O}}, \underline{K}. The
% symmetric groups of pages 17-18 are \mathfrak{S}_p, and A(\mathcal{X}, p) is
% his abbreviation of A(\mathcal{X}, \mathfrak{S}_p).
%
% Readings flagged and not settled. On page 1 the struck third line of the
% wrapper, of which only "de R.R." carries. On page 2 the word after E_0 (a
% noun beginning in g and ending in -ti., abbreviated), the prefix inside the
% quotation marks in "diviseurs « f-équivalents » à 0", and the exact wording
% of the proposition's second family of generators. On page 3 the first word of
% the page. On page 4 the phrase after "faisceau loc. libre de rang". On pages
% 6 and 7 most of the connective prose, which is a fast palimpsest; page 7 is
% crossed by long diagonal cancellation strokes from its middle down. On page 9
% the brace under the k-th term. On page 11 the sentence that closes the page.
% On page 13 the right-hand side of the last displayed line, read 0 and which
% the argument would want to be 1.
\documentclass[11pt,a4paper]{article}
% Le préambule commun à toutes les transcriptions du fonds Grothendieck.
%
% Il tient en une idée : **l'incertitude se compose, elle ne se résout pas.**
% Une transcription qui lisse un mot douteux a détruit la seule chose pour
% laquelle on la faisait — un lecteur doit pouvoir savoir, à chaque endroit,
% ce qui était sur le feuillet et ce qui a été ajouté. Les six macros qui
% suivent sont donc l'appareil critique, et non de la décoration.
%
% Les mêmes commandes sont reconnues par `scripts/render.mjs`, qui produit la
% vue de lecture du site. Une macro ajoutée ici doit l'être aussi là-bas,
% faute de quoi le rendu échouera bruyamment — ce qui est voulu : un rendu
% silencieusement faux est pire qu'un rendu absent.

\ProvidesPackage{grothendieck}[2026/08/08 Transcriptions du fonds Grothendieck]

\usepackage{amsmath,amssymb,amsthm}
\usepackage{mathrsfs}
\usepackage{stmaryrd}
% Les diagrammes commutatifs sont partout dans ce fonds : c'est la langue
% même de la géométrie algébrique de Grothendieck, et une transcription qui
% les rendrait en prose ne transcrirait rien.
\usepackage{tikz-cd}
% Français par défaut — c'est la langue des feuillets — mais l'anglais est
% chargé aussi : la traduction et le résumé sont des documents anglais, et la
% typographie française (espaces avant les deux-points) y serait une faute.
% Ces documents basculent par \selectlanguage{english} en tête de corps.
\usepackage[english,french]{babel}
\usepackage{fontspec}
\usepackage{geometry}
\geometry{a4paper,margin=25mm,marginparwidth=28mm}
\usepackage{marginnote}
\usepackage{xcolor}
% ulem, not soul. `soul`'s \st and \ul are built on a character-by-character
% scan that XeLaTeX's font machinery breaks: the document fails with an
% undefined control sequence rather than degrading. ulem does the same two jobs
% and is Unicode-engine safe. `normalem` stops it hijacking \emph, which the
% transcriptions use for Grothendieck's own emphasis.
\usepackage[normalem]{ulem}
\usepackage{hyperref}
\hypersetup{colorlinks=true,linkcolor=black,urlcolor=black,pdfborder={0 0 0}}

% --- Métadonnées du lot -----------------------------------------------------
% Reprises telles quelles par le script de rendu, qui les affiche en tête de
% la vue de lecture. Elles fixent ce que le fichier prétend couvrir : sans
% elles, une transcription flotte sans point d'attache dans un fonds de seize
% mille feuillets.

\newcommand{\folder}[1]{\gdef\grothFolder{#1}}
\newcommand{\batch}[1]{\gdef\grothBatch{#1}}
\newcommand{\leaves}[2]{\gdef\grothFirst{#1}\gdef\grothLast{#2}}
\newcommand{\foldertitle}[1]{\gdef\grothFoldertitle{#1}}
\gdef\grothFolder{?}\gdef\grothBatch{?}\gdef\grothFirst{?}\gdef\grothLast{?}\gdef\grothFoldertitle{}

% --- Le résumé --------------------------------------------------------------
%
% Il ouvre la lecture modernisée, sous le titre « Résumé », et il est composé
% en petit. Ce n'est pas un ornement : c'est le seul passage du document qui
% ne soit pas des mathématiques, et un lecteur doit pouvoir le reconnaître
% avant de l'avoir lu — puis le sauter s'il connaît déjà le sujet, ou s'y
% arrêter s'il ne le connaît pas. Le filet qui le suit marque la reprise du
% corps.

\newenvironment{resume}
  {\small\setlength{\parskip}{0.45em}}
  {\par\medskip\hrule\medskip}

% --- Mots-clés ---------------------------------------------------------------
%
% En anglais, en clôture du résumé de la lecture modernisée : le vocabulaire
% moderne sous lequel chercher ce que les feuillets construisent. Le manifeste
% du site (`npm run manifest`) les extrait pour étiqueter la cote entière —
% cette ligne est la source unique des tags, il n'y a pas de fichier à côté.
\newcommand{\keywords}[1]{\par\smallskip\noindent
  {\footnotesize\textsc{Keywords} --- \textit{#1}}\par}

% --- L'appareil critique ----------------------------------------------------

% Le numéro de feuillet, en marge. C'est l'ancre qui permet de régler un
% désaccord sur une formule en regardant *un* feuillet plutôt qu'une chemise
% de six cents. Le site s'en sert aussi pour tourner les pages du fac-similé
% quand on fait défiler la transcription.
\newcommand{\leaf}[1]{%
  \marginnote{\footnotesize\color{gray!70}\textsf{#1}}%
  \label{leaf:#1}\ignorespaces}

% Illisible : on ne devine pas. Un mot inventé qui se lit comme les autres est
% le pire résultat possible.
\newcommand{\ill}{\textcolor{red!60!black}{[\,\ldots\,]}}

% Lecture douteuse : le mot est donné, et signalé comme incertain.
\newcommand{\uncertain}[1]{\uline{#1}}

% Ajout de l'éditeur — une lettre manquante, un mot sous-entendu.
\newcommand{\add}[1]{\textcolor{blue!50!black}{[#1]}}

% Biffé par Grothendieck lui-même. Ce qu'il a rayé fait partie du document :
% une rature dit souvent où la pensée a changé de direction.
\newcommand{\struck}[1]{\sout{#1}}

% Note du transcripteur, sur l'état matériel du feuillet.
\newcommand{\note}[1]{\par\smallskip{\small\color{gray!60!black}\textit{[#1]}}\par\smallskip}

% Note marginale de Grothendieck, distincte du corps du texte.
\newcommand{\marginal}[1]{\par\smallskip\noindent\hspace{1em}%
  {\small\color{gray!60!black}#1}\par\smallskip}

% --- Titre ------------------------------------------------------------------

\AtBeginDocument{%
  \begin{center}
    {\large\bfseries Fonds Alexandre Grothendieck --- cote n\textsuperscript{o}~\grothFolder}\\[2pt]
    {\itshape\grothFoldertitle}\\[4pt]
    {\small feuillets \grothFirst--\grothLast\ (lot \grothBatch)}\\[2pt]
    {\footnotesize\color{gray!60!black}Transcription établie d'après le fac-similé numérisé
     par l'Université de Montpellier.\\
     Les lectures incertaines sont soulignées, les illisibles notées [\,\ldots\,],
     les ajouts entre crochets.}
  \end{center}
  \vspace{1em}}

\endinput


\folder{2}
\batch{1}
\pages{1}{20}
\dating{[1957]}
\watermark{Édition de démonstration}
\foldertitle{Théorème de Riemann-Roch / RR [Riemann-Roch] Notes Antiques (antérieur[es à] 1958) : notes manuscrites (s.d.).}

\begin{document}

\section*{\uncertain{Classique pour les Courbes}}

\note{Le titre est de lui : il est encadré au crayon sur la chemise qui ouvre
le lot (page 1), laquelle ne porte rien d'autre. Au-dessus de l'encadré, biffé
de longs traits obliques : \struck{Fonctorialité des variances} \struck{\ill{}
de R.R.}}

\page{2}
$X$ courbe non singulière.

$E$ anneau des classes de faisceaux cohérents mod. les $(F) - (F') - (F'')$,
\marginal{$\rightarrow$ \uncertain{suite} exacte $0 \to F' \to F \to F'' \to 0$}
où \ill{}

$E_{0}$ \uncertain{noyau} de l'\ill{}\note{le mot, abrégé, n'a pas été lu ; il
commence par un $g$ et se termine en \texttt{-ti.}} \quad (\ill{})

$\mathcal{D}$ sous-anneau engendré par les faisceaux \uncertain{définis} du
type $k_{x}$ ; \quad $\mathcal{D}$ = groupe des diviseurs. \quad ${}^{*}$C'est
l'annulateur de $E_{0}$, donc un idéal.

$N$ sous-groupe de $E$ engendré par les $(F) - (F') - (F'')$.

$\mathcal{D} \cap N$ \struck{\ill{}} diviseurs « \uncertain{$f$}-équivalents »
à $0$, \struck{identiques à}\note{le préfixe entre guillemets est une haste
bouclée ; il pourrait aussi bien être un $l$ pour « linéairement », ou le
signe $\int$ que la page 3 emploie pour la caractéristique.}

$\mathcal{L}$ groupe des diviseurs équivalents à $0$.

$\Phi$ sous-algèbre de $E$ engendrée par les \struck{faisceaux} loc. libres
\quad $\simeq$ algèbre des classes de fibrés vectoriels, \ill{}

$\Phi_{N}$ (quotient de $\Phi$ par le sous-groupe ($=$ idéal) engendré par les
$c(E)$, $c(E')$, \ill{})


${}^{*}$ $E = \Phi + \mathcal{D}$, \quad $\Phi \cap \mathcal{D} = 0$

${}^{*}$ $E/N \simeq \Phi/\Phi_{N} \simeq \mathbb{Z} + \mathfrak{P}$ \quad (1)

${}^{*}$ \struck{\ill{}}


(1) $\mathbb{Z} + \mathfrak{P} \to E/N$ fait correspondre à $(1,0)$ la classe
de $\underline{\mathcal{O}}_{X}$, et sur $\mathfrak{P}$ \ill{} déduit par
passage au quotient de l'injection $\mathcal{D} \subset E$.


\textbf{Proposition.} $N$ est engendré \add{par} les $c(E) - c(E') - c(E'')$,
où $E' \to E \to E'' \to 0$ est une suite exacte de fibrés vectoriels, et les
\[
  c(E) - c(F) - c\bigl(\underline{\mathcal{O}}_{X}(E)/f\,\underline{\mathcal{O}}_{X}(F)\bigr),
\]
où $E$ et $F$ sont des fibrés linéaires et $f : F \to E$ un \ill{}, i.e.
$f \in H^{0}(X, \underline{\mathcal{O}}_{X}(F^{*} \otimes E))$.

\add{En} d'autres \add{termes}, \ill{} les $\lambda(D + D') - \lambda(D') - \ill{}$,
où $D$ et $D'$ \uncertain{diviseurs}, et où $\lambda(\Delta)$ est la classe du
faisceau inversible \add{ou} fibré vectoriel associé au diviseur $\Delta$. Ce
qui \uncertain{se ramène} à $D \geqslant 0$.

\page{3}
\ill{}\note{un mot de huit ou neuf lettres, commençant par L-i et finissant
en \texttt{-tion}} précédant la proposition est équivalente aux suivantes :
\[
  \mathcal{D} + N = E_{0}, \qquad \mathcal{D} \cap N = \mathcal{L},
  \qquad \struck{\Phi + N = E}, \qquad \Phi \cap N = \Phi_{N} .
\]


\textbf{Définition de l'application} $E \longrightarrow \mathbb{Z} + \mathfrak{P}$ :

\begin{itemize}
  \item sur $\mathcal{D}$ : coïncide avec l'application canonique
        $\mathcal{D} \to \mathfrak{P}$ ;
  \item sur $\Phi$ : soit $E$ un fibré vectoriel, \uncertain{fibre} de type
        $k^{n}$ et de classe de Chern $c^{1}(E) \in \mathfrak{P}$ ; on prend
        $h(E) = (n,\, c^{1}(E))$.
\end{itemize}

Soit $F$ un faisceau algébrique cohérent, $c(F) \in \mathbb{Z} + \mathfrak{P}$ ;
on voit que
\begin{equation*}
  \int(F) \;=\; K_{1}\!\left( c(F) \cdot \left(1 - \tfrac{1}{2}K\right) \right)
  \qquad (K \in \mathfrak{P}) \tag{1}
\end{equation*}
où $K$ est la classe canonique ; $K_{1}$ est la forme linéaire sur
$\mathbb{Z} + \mathfrak{P}$ qui est $0$ sur $\mathbb{Z}$, et sur $\mathfrak{P}$
l'application « degré » $\mathfrak{P} \to \mathbb{Z}$. Pour vérifier (1), il
suffit de le vérifier pour $F = k_{x}$ et pour $F = \underline{\mathcal{O}}_{X}$.

\note{Le symbole $\int$ est, sur le feuillet, une haste bouclée en forme de
$s$ long ; c'est la caractéristique d'Euler-Poincaré que la formule calcule.
Le symbole $\mathfrak{P}$ est un capitale cursive barrée d'un trait horizontal,
distincte du $\mathcal{L}$ de la page 2.}


\begin{tikzcd}[column sep=small, row sep=small, nodes={font=\scriptsize}]
\mathfrak{P} & \mathcal{D} \arrow[rd] & & \\
\mathcal{L} = \mathcal{D} \cap N \arrow[ru] \arrow[rd, "?"'] & & \mathcal{D} + N = E_{0} \arrow[r] & E \\
& N \arrow[ru] & &
\end{tikzcd}

\note{Sur le feuillet, $\mathfrak{P}$ surmonte $\mathcal{D}$ ; deux accolades
embrassent la partie droite du dessin, celle du haut portant $\Phi$, celle du
bas $\mathbb{Z}$, et sous celle du bas : « $E/N$ classes de faisceaux ».}

\page{4}
$X$ \struck{préschéma} [noethérien] admettant \struck{aussi} un Module
inversible ample, $\dim X \leqslant 1$.

$K^{\bullet}(X)$ augmenté vers $H^{0}(X, \mathbb{Z}) = \mathbb{Z}^{\pi_{0}(X)}$,
$\quad E \rightsquigarrow (\mathrm{rg}_{i}(E))_{i \in \pi_{0}(X)}$.

L'idéal d'augmentation $I(X)$, \ill{} on a aussi
\[
  \left.
  \begin{array}{l}
    \text{si } x \in I(X),\ \ill{} \\[2pt]
    \gamma^{i}(x) = 0 \quad \text{pour } i \geqslant 2
  \end{array}
  \right\}
  \qquad F^{i}_{\gamma}(K^{\bullet}(X)) = 0 \ \text{si } i \geqslant 2 .
\]

Donc
\[
  K^{\bullet}(X) \simeq \mathbb{Z}^{\pi_{0}(X)} \times K^{(1)}(X),
  \qquad
  K^{(1)}(X) = \mathrm{Pic}(X) \simeq \prod_{i \in \pi_{0}(X)} \mathrm{Pic}(X_{i}),
\]
\struck{Donc $E$} $K^{\bullet}(X)$ \uncertain{entièrement déterminé} \ill{}
comme $\lambda$-anneau augmenté ; \ill{} par les $\mathrm{Pic}(X_{i})$,
$i \in \pi_{0}(X)$.

Si $E$ est un \struck{fibré vectoriel} faisceau loc. libre de rang
\uncertain{$0$} \ill{} sur $X$, \ill{}
\[
  \mathrm{cl}(E) = \bigl[\mathrm{rg}(E),\, c^{1}(E)\bigr]
  \qquad \text{où } c^{1}(E) = \det(E),
\]
\[
  \Bigl[\ \text{i.e.}\quad
  \mathrm{cl}(E) = \sum_{i} \mathrm{rg}_{i}(E)\cdot 1_{i} + \bigl[\mathrm{cl}(\det E) - 1\bigr]
  = \sum_{i} \bigl(\mathrm{rg}_{i}(E) - 1\bigr) 1_{i} + \mathrm{cl}(\det E)\ \Bigr]
\]
\note{la ligne d'au-dessus porte, insérée, $(\mathrm{rg}_{i}(E))_{i \in I}$.}

On a $B^{\bullet}(X) = \mathbb{Z}^{\pi_{0}(X)} \times \mathrm{Pic}(X)$, et
$K^{\bullet}(X) \xrightarrow{\ c\ } B^{\bullet}(X)$ est l'identité, avec les
identifications précédentes.

$K_{\bullet}(X)$ augmenté vers $\mathbb{Z}^{\pi_{0}^{(1)}(X)}$,
$\quad F \rightsquigarrow (\mathrm{lg}_{i}(F))_{i \in \pi_{0}^{(1)}(X)}$
\quad [$\pi_{0}^{(1)}(X)$ = ens. des composantes irréductibles de $\dim 1$].

Le noyau est $K_{0}(X)$ ; il est isomorphe à
\[
  \mathbb{Z}^{k^{(0)}(X)} \times \mathrm{Pic}'(X)
  \;=\; \mathbb{Z}^{k^{(0)}(X)} \times \prod_{i \in \pi_{0}^{(1)}(X)} K_{0}(X_{i})
\]
\note{sous l'indice du produit il écrit « \ill{} des composantes connexes de
$\dim 1$ », et il souligne « $\dim 1$ » deux fois.}

\page{6}
Pour écrire les opérations de \uncertain{variance} entre $K^{\bullet}(X)$ et
$K_{\bullet}(X)$, \ill{} (\ill{} et \ill{} $1$ de $K_{\bullet}(X)$ \ill{} pour
simplifier l'écriture.) \quad \emph{$X$ courbe.}


(a $1^{\circ}$) \quad $\dim X = 0$.

Les opérations de $K^{\bullet}(X)$ \ill{}. \quad $1 \in K_{\bullet}(X)$ donné
par \ill{} :
\[
  \mathrm{lg}(1) = (\mu_{i}) \in \mathbb{Z}
  \qquad \text{seul invariant de la situation}
\]
\struck{\ill{}} \quad ici $c^{\bullet} = \mathrm{ch}^{\bullet} = \mathrm{id}_{K}$ ?
\[
  K^{\bullet}(X) \xrightarrow{\ \mathrm{rg}\ } \mathbb{Z},
  \qquad
  K_{\bullet}(X) \xrightarrow{\ \mathrm{lg}\ } \mathbb{Z}
\]
(isomorphismes de $\lambda$-anneaux), et $K^{\bullet}(X) \rightsquigarrow
K_{\bullet}(X)$ : \uncertain{pour} $2$ choix !


$2^{\circ}$) \quad $\dim X = 1$. \quad \ill{} on a les invariants suivants :

\begin{enumerate}
  \item[a)] L'ens. fini $K = \pi_{0}(X)$ des comp. irréd. de $X$ ;
  \item[b)] les \ill{} ; \quad $\delta \in K_{0}(X)$ \ill{} ; \quad
        $(\mu_{i})_{i \in K}$ (dans \ill{} $\mathrm{lg}_{i}(\underline{\mathcal{O}}_{X})$),
        $\bigl(= 1 - \sum \mu_{i} 1_{i}\bigr)$ \note{les deux lignes de ce
        point sont réunies par une accolade sur le feuillet} ;
  \item[c)] Le groupe \quad $K^{(1)} = \mathrm{Pic}(X)$ ;
  \item[d)] Le groupe \quad $K^{\varepsilon}_{(0)} = K_{(0)}(X_{i})$,
        \struck{\ill{}} $\prod_{i \in K} K^{i}_{(0)}$ ;
  \item[d)] D'autres \ill{} : $u_{i} : K^{(1)} \to K^{\varepsilon}_{(0)}$,
        \struck{\ill{}} $i \in K$, \quad
        $L \rightsquigarrow \mathrm{cl}\bigl((L-1)\underline{\mathcal{O}}_{X_{i}}\bigr)$.
\end{enumerate}

\marginal{un invariant, \uncertain{dépendant} du \ill{} \uncertain{choix} \ill{}
la classe \ill{} \add{de} l'élément $1$ de $K_{\bullet}(X)$}

\ill{} on écrit \ill{} comment
\[
  K^{\bullet}(X) = \mathbb{Z} \times K^{(1)}
  \xrightarrow[\ = \mathrm{id}\ ]{\ c^{\bullet} = \mathrm{ch}^{\bullet}\ }
  B^{\bullet}(X) \simeq \mathbb{Z} \times K^{(1)} ,
\]
\[
  K_{\bullet}(X) = \mathbb{Z}^{K} \times K_{(0)}
  \ \struck{\ill{}} \ = \ B_{\bullet}(X) \simeq \mathbb{Z}^{K} \times K_{(0)} .
\]
$K^{\bullet}(X) \rightsquigarrow K_{\bullet}(X)$ : il suffit \uncertain{pour}
\ill{} définir des $K_{\bullet}(X)$ \ill{} opérations \ill{} de degré filtrant
$\geqslant 1$ \ [\ill{}], \ill{}
\[
  \mathbb{Z}^{K} \longrightarrow K_{(0)},
  \qquad \struck{\ill{}} \qquad
  \xi \rightsquigarrow \bigl(\uncertain{\lambda_{i}}\, u_{i}(\xi)\bigr)_{i \in K} .
\]

\page{7}
\uncertain{Homo.} de $K^{\bullet}(X)$ dans $K_{\bullet}(X)$, défini par
$1_{\bullet} \in K_{\bullet}(X)$ défini par
\[
  1_{\bullet} = \sum \mu_{i}\, 1_{i} + \delta_{\varepsilon} .
\]
Alors \ill{}
\[
  u(n, \xi) \;=\; \sum n\,\mu_{i}\, 1_{i}
  \;+\; \Bigl( n\,\delta_{\varepsilon} + \sum \mu_{i}\, u_{i}(\xi) \Bigr),
  \qquad n \in \mathbb{Z},\ \xi \in K^{(1)} .
\]

\struck{\ill{} on définit \ill{} de \ill{} même
$K_{\bullet}(X) \xrightarrow{\ c_{\bullet} = \mathrm{ch}_{\bullet}\ } B_{\bullet}(X)$,
$c_{\bullet}(\xi_{\bullet}) = \ill{}$}\note{le bas du feuillet est barré de
longues obliques ; la formule qu'elles traversent fait intervenir les
$\mathrm{lg}_{i}(\xi_{i})$ et une soustraction de $\sum \ill{}\, 1_{i}$, et
n'a pas été lue.}

\section*{Contre-exemples}

\note{Le mot est de lui : il est seul, en haut à droite de la chemise qui ouvre
le second lot (page 8).}

\page{9}
\[
  X \xrightarrow{\ i\ } X \times \mathbb{P}^{r},
  \qquad F \otimes L .
\]
\[
  \underline{K}(X \times \mathbb{P}^{r})
  = \underline{K}(X) \otimes \underline{K}(\mathbb{P}^{r}) \big/ \bigl(\xi = 1 - L^{-1}\bigr) .
\]
\[
  i_{!}(x) = p^{!}(x) \otimes q^{!}(\xi^{r}) = x \otimes \xi^{r},
\]
\[
  \widetilde{C}(i_{!}(x)) = \widetilde{C}(x \otimes \xi^{r})
  = \widetilde{C}(x) * \widetilde{C}(\xi^{r}) .
\]

$\underline{K}(X) \to \underline{K}(X \times \mathbb{P}^{r})$ est
\uncertain{injectif} ; on a le carré

\begin{tikzcd}[column sep=small, row sep=small, nodes={font=\scriptsize}]
\underline{K}(X) \arrow[r, "\text{inj.}"] \arrow[d, "\widetilde{c}_{X}"'] & \underline{K}(X \times \mathbb{P}^{r}) \arrow[d, "\widetilde{c}_{X \times \mathbb{P}^{r}}"] \\
\widetilde{A}(X) \arrow[r] & \widetilde{A}(X \times \mathbb{P}^{r})
\end{tikzcd}

\ill{} \quad $p_{*} i_{*}(x) = (p i)_{*}\, x = x$.

\[
  \xi^{r} = \sum_{i=0}^{r} \binom{r}{i} (-1)^{i} L^{-i},
  \qquad
  \widetilde{C}(\xi^{r}) = \prod_{i=0}^{r} (1 - i\eta)^{(-1)^{i} \binom{r}{i}} .
\]
\[
  \widetilde{C}(i_{!}(x))
  = \prod_{i=0}^{r} \Bigl[\, \widetilde{C}(x) * (1 - i\eta) \,\Bigr]^{(-1)^{i} \binom{r}{i}}
\]
\[
  \struck{= \prod_{i=0}^{r} \Bigl[ (1 - i\xi)^{\varepsilon(x)} \sum_{i=0}^{\infty} C^{i}(x) \Bigr]}
\]

\[
  \widetilde{C}(x) * (1 - i\eta)
  = \widetilde{C}(x) * \bigl[1,\ 1 - i\eta\bigr] - \widetilde{C}(x)
  = \Bigl[\ \varepsilon(x),\ 1 + \cdots + \sum_{j=0}^{k} C^{j}(x)\,(-i\eta)^{k-j}\ \Bigr]
\]
\note{une accolade court sous le dernier terme et porte « $k$\ill{} terme
\ill{} $C^{k}$ de la seconde \ill{} : $\widetilde{C}(x) * [1,\, 1 - i\eta]$ » ;
le milieu de cette glose n'a pas été lu.}

\[
  \widetilde{C}(\xi^{r}) = 1 + (-1)^{r+1}(r-1)!\ \eta^{r},
\]
d'où
\[
  \widetilde{C}(i_{!}(x)) = \widetilde{C}(x) * \Bigl(1 + (-1)^{r+1}(r-1)!\ \eta^{r}\Bigr)
\]
\[
  = \Bigl(1 + (-1)^{r+1}(r-1)!\ \eta^{r}\Bigr)^{\varepsilon(x)}
    \Bigl( C(x) * \bigl(1 + (-1)^{r+1}(r-1)!\ \eta^{r}\bigr) \Bigr)
\]
\[
  = 1 + \sum_{i=1}^{\infty} Q_{i}\bigl(c^{1}(x), \ldots, c^{i}(x);\ 0, \ldots,
    \uncertain{\alpha^{(i/r)}}, 0, \ldots, 0\bigr)
\]
\note{les exposants $(-1)^{r+1}$ de ce bloc sont écrits $(-1)^{r-1}$ à la page
suivante ; la page ne tranche pas.}

\marginal{$b_{\gamma} = 1 - L(-\gamma)$, \quad $1 - b_{\gamma} = L(-\gamma)$}

\page{11}
Or $Q_{i}$ \struck{\ill{}} \ill{} en \ill{} contenant \ill{} \uncertain{singulier} ;
en dehors, \uncertain{sachant que} $(\eta^{r})^{n} = 0$ si $n > 1$. Donc on
trouve
\[
  Q_{i}(\cdots) =
  \begin{cases}
    0 & \text{si } i < r, \\[4pt]
    (-1)^{r-1}(r-1)!\ Q^{(r)}_{i-r}\bigl(c^{1}(x), \ldots, c^{i}(x)\bigr) &
  \end{cases}
\]
où $Q^{(r)}_{i-r}\bigl(c^{1}(x), \ldots, c^{i}(x)\bigr)$ \ill{} est le
coefficient de \ill{} (\ill{} de poids $i - r$) \ $d^{r}$ dans
$Q_{i}(c^{1}, \ldots, c^{i};\, d^{1}, \ldots, d^{i})$ \quad ($r \leqslant i$)
\quad (\ill{} Attention \ill{}).

\[
  \widetilde{C}(i_{!}(x)) = 1 + (-1)^{r-1}(r-1)!\ \eta^{r}
  \Bigl[\ \varepsilon(x) + \sum_{i=0}^{\infty} Q^{(r)}_{i}\bigl(c^{1}(x), \ldots, c^{i}(x)\bigr) \Bigr]
\]

\emph{Comparons avec R.R.} \quad
$\widetilde{C}(i_{!}(x)) = \struck{\ill{}}\ (i_{*})_{r}\bigl(\widetilde{C}(x)\bigr)$.
On voit que
\[
  (i_{*})_{r}\Bigl[\ \eta \sum_{0}^{\infty} x^{i} \Bigr]
  = 1 + (-1)^{r-1}(r-1)!\ \eta^{r}
    \Bigl[\ 1 + \sum_{i=0}^{\infty} Q^{(r)}_{i}\bigl(x^{i}, \ldots, x^{i}\bigr) \Bigr]
\]

N.B. $Q^{(r)}_{0}(c^{1}, \ldots, c^{r}) = 0$, i.e.
$Q_{r}^{\sharp}(c^{1} \ldots c^{r};\, 0, \ldots, 0, d^{r}) = 0$, i.e.
$(1 + c^{1} \ill{} - c^{r}) * (1 + d^{r})$ \ill{} à filtration $\geqslant r+1$.

De \ill{} on voit
\[
  Q^{(r)}_{1}(c^{1}) - \struck{\ill{}}
  = -\ \frac{(r+1-1)!}{(r-1)!\,(1-1)!}\ c^{1} = -\,r\,c^{1} .
\]

\ill{} les coefficients numériques \ill{} le $(r-1)!$ \ill{} provenant des
$Q^{(r)}_{i}$.

\page{13}
\uncertain{Conséquence.} Soit dans $Y$ un \uncertain{hyperplan}, \ill{}
$x' \in A^{1}(X)$ sa classe ; on a
\[
  \widetilde{C}_{X}(Y) = \bigl(1 - \struck{c^{1}(x)}\ x'\bigr)^{-1} = \sum x'^{\,n} ,
\]
\[
  \widetilde{C}_{X \times \mathbb{P}^{r}}\bigl(i_{!}(\underline{\mathcal{O}}_{Y})\bigr)
  = \widetilde{C}_{X \times \mathbb{P}^{r}}(Y \times \mathbb{P}^{r})
  = 1 + (-1)^{r-1}(r-1)! \sum_{i=1}^{\infty}
      Q^{(r)}_{i}\bigl(x',\, (x')^{2}, \ldots, (x')^{i}\bigr)
\]
\[
  = 1 + (-1)^{r-1}(r-1)!\ \eta^{r} \sum_{i=1}^{\infty}
      Q^{(r)}_{i}\bigl(x',\, (x')^{2}, \ldots, (x')^{i}\bigr)
\]
\note{une accolade sous la somme porte $c_{i,r}\,(x')^{r}$.}

Si donc $(r-1)!\,x' = 0$, \ill{} \quad
$\widetilde{C}_{X \times \mathbb{P}^{r}}(Y \times \mathbb{P}^{r}) = \uncertain{0}$,
donc \ill{} \uncertain{pour prendre} \ill{} $Y'$ \ill{} \uncertain{soit}
$x' = 0$.

\note{Le dernier membre est lu $0$ ; l'argument attendrait $1$, et la page ne
tranche pas.}

\section*{Modules sur un topos annelé où un groupe opère}

\note{Le titre est de l'éditeur : ces quatre feuillets ne sont précédés
d'aucune chemise inscrite dans ce lot, et leur sujet n'est pas celui des
« Contre-exemples » qui précèdent.}

\page{15}
\struck{\ill{}} $\mathcal{X}$ Topos annelé, $\mathcal{O}_{\mathcal{X}} =
\mathcal{O}$ son anneau (supposé commutatif).

Pour tout groupe $G$, on considère la catégorie
\[
  A(\mathcal{X}, G)
\]
des Modules sur $\mathcal{X}$ où $G$ opère. On a, pour un \add{homomorphisme}
$u : G \to G'$, un homomorphisme de restriction
\[
  u^{*} : A(\mathcal{X}, G') \longrightarrow A(\mathcal{X}, G)
\]
et \ill{} un homomorphisme \ill{}
\[
  u_{*} : A(\mathcal{X}, G) \longrightarrow A(\mathcal{X}, G') .
\]
Lorsque $u$ est injectif, \ill{} on a \struck{\ill{}} un isomorphisme
canonique \ill{}
\[
  u^{*} u_{*}(F) \simeq F^{(G'/G)} \quad \ill{}
\]
Si on se borne aux sous-catégories \uncertain{pleines} \ill{} formées des
Modules \uncertain{loc.} libres, resp. libres de type fini, resp. loc. libres
de type fini) sur $\mathcal{O}[G]$, elles sont stables par les foncteurs
précédents. Elles \ill{} par les foncteurs que \ill{} définit maintenant.

Dans $A(\mathcal{X}, G)$ on a l'opération $\otimes$ (commutative et
associative), \ill{}

\page{16}
\ill{} on a les isomorphismes fonctoriels
\[
  u^{*}(F' \otimes G') \simeq u^{*}(F') \otimes u^{*}(G')
\]
et pour $u_{*}$ \ill{} la « formule de projection »
\[
  u_{*}\bigl(F \otimes u^{*}(G')\bigr) \simeq u_{*}(F) \otimes G' .
\]

[On a aussi les opérations $\lambda^{i}$ (et \ill{} les opérations
tensorielles) qui \ill{}. Mais \uncertain{il est désirable} \ill{} plus bas des
opérations \ill{} $\otimes$]

\ill{} on \uncertain{définit} \ill{} le produit tensoriel « extérieur » :
\[
  F \in \mathrm{Ob}\ A(\mathcal{X}, G), \qquad G \in \mathrm{Ob}\ A(\mathcal{X}, H),
\]
\[
  F \boxtimes G = \mathrm{pr}_{1}^{*}(F) \otimes \mathrm{pr}_{2}^{*}(G)
  \ \in\ \mathrm{Ob}\ A(\mathcal{X}, H \times G) .
\]
\ill{} les \ill{} inverses \ill{}
\[
  u^{*}(F \boxtimes G) \simeq u^{*}(F) \boxtimes u^{*}(G),
\]
\[
  F \otimes G \ \rightsquigarrow \ = (\mathrm{diag}_{G})^{*}(F \boxtimes G) .
\]

\page{17}
On désigne par
\[
  u_{pq} : \mathfrak{S}_{p} \times \mathfrak{S}_{q} \longrightarrow \mathfrak{S}_{p+q}
\]
l'injection canonique, et on pose, pour
$F \in \mathrm{Ob}\ A(\mathcal{X}, p)$, $G \in \mathrm{Ob}\ A(\mathcal{X}, q)$ :
\[
  F * G = u_{pq\,*}\bigl(F \boxtimes G\bigr) \ \in\ \mathrm{Ob}\ A(\mathcal{X}, p+q) .
\]
L'opération $F * G$ \ill{} est « commutative » et « ass. ». \quad \ill{}
\struck{\ill{}}

\note{$A(\mathcal{X}, p)$ abrège $A(\mathcal{X}, \mathfrak{S}_{p})$ ; son
$\mathfrak{S}$ est un $S$ bouclé.}

\page{18}
Formules \quad $\bigl[F, G \in \mathrm{Ob}\ A(\mathcal{X})\bigr]$ :
\[
  T^{n}(F + G) \simeq \coprod_{p+q=n} T^{p}(F) * T^{q}(G),
\]
\[
  T^{m}\bigl(T^{n}(F)\bigr) \simeq w_{m,n}^{*}\bigl(T^{mn}(F)\bigr),
\]
\[
  T^{m}(F \otimes G) \simeq T^{m}(F) \otimes T^{m}(G)
\]
\[
  \bigl[\ w_{m,n} : \mathfrak{S}_{m} \longrightarrow \mathfrak{S}_{\ill{}}
  \ \ill{} \text{convenable} \ \bigr]
\]

Nous \ill{} les Modules $F$ ayant la \ill{} \ill{} $\mathcal{O}_{\mathcal{X}}(\mathfrak{S}_{n})$
\ill{} la propriété suivante :
($*$) \quad $F \rightsquigarrow$ le \ill{} \ill{} les faisceaux
\[
  u_{pq\,*}\bigl(F_{p_{1} p_{2} \ldots p_{r}}\bigr)
  \qquad \Bigl(\sum p_{i} = n,\ \ill{},\ p_{i} \geqslant 0,\ r \in \mathbb{N}\Bigr)
\]
\ill{} des \emph{faisceaux} \ill{} loc. \ill{}
$\bigl(\mathfrak{S}_{p_{1}} \times \mathfrak{S}_{p_{2}} \times \cdots
\times \mathfrak{S}_{p_{r}}\bigr)$.

\section*{Cas des schémas}

\note{Les mots sont de lui : la chemise qui ouvre le troisième lot (page 19)
porte « Cas des schémas », souligné, puis le sommaire de ce qui suit —
« Filtration des $\lambda$-anneaux, classes de Chern universelles ;
$\lambda$-anneaux "spéciaux" i.e. avec éléments "géométriques" ;
Application », suivi d'un signe de validation.}

\page{20}
$K$ $\lambda$-anneau \struck{\ill{}} \emph{spécial},
$(x_{s})_{s \in S}$ \struck{\ill{}} \emph{famille} de générateurs (pour
la structure) \ill{}

On \uncertain{introduit} \ill{} $\lambda$\ill{} \quad
$K \xrightarrow{\ \rho\ } \widetilde{A}$ \quad ($\widetilde{A}$, \ill{}).

Il n'y a \ill{} plus d'applications universelles, et \ill{} \uncertain{avec}
l'aide des $x_{s}$, on \ill{}

E Ne \ill{} que \ill{} \ill{} structure $S$ \ill{} $+$, $*$ et \ill{}

\emph{Exemple} : \uncertain{considérons} \struck{\ill{}} $B$ le \ill{}
spécial \ill{} engendré par \ill{} générateurs $X_{s}$ ($s \in S$) soumis à
\struck{\ill{}} ($X_{s}$ \ill{} $\varepsilon(s)$), i.e.
\[
  \lambda^{i}(X_{s}) = 0 \ \text{si } i > \varepsilon(s),
  \qquad \lambda^{\varepsilon(s)}(X_{s}) \ \text{inversible} .
\]
D'où
\[
  B = \mathbb{Z}\Bigl[\bigl(\lambda^{i}(X_{s})\bigr)_{s \in S,\ 1 \leqslant i \leqslant \varepsilon(s)}\Bigr]
      \Bigl[\bigl(\lambda^{\varepsilon(s)}(X_{s})\bigr)^{-1}\Bigr] .
\]
Prenons $y = $ la famille $(X_{s})_{s \in S}$. \ill{}
\[
  CB = \mathbb{Z}\Bigl[\bigl(C^{i}_{s}\bigr)_{\struck{\ill{}}}\Bigr],
  \qquad \ill{} \quad
  \widetilde{C}(X_{s}) = \Bigl[\varepsilon(s),\ 1 + C^{1}_{s} \cdots \ill{}\Bigr] .
\]
Il y a sur $\widetilde{C}B$ une filtration, d'\ill{}

\ill{} $\gamma$-filtration \ill{} $\widetilde{C}K$, \quad $GK = \mathrm{gr}\,K$,
\quad $\psi_{K} : GK \to G\widetilde{C}K \simeq CK$ \ill{}

\textbf{Proposition 1.} \ill{}

\textbf{Proposition 2.} \ill{} \quad
$\varphi_{K} : GK \longleftarrow CK$, \quad
$\varphi_{K}\psi_{K} = (-1)^{i}\,(i-1)! \cdot \mathrm{id}$ \ill{}

\note{Le bas de ce feuillet est une suite serrée d'énoncés dont seuls les
symboles portent ; la prose qui les relie n'a pas été lue.}

\end{document}
